% Fichier GÉNÉRÉ (v0.23.3) — les comptes EXACTS des journaux
% récupérés, la source unique de vérité partagée par les cinq pièces.
% (le worklog et le DOI recalculés par
% scripts/vague15/regenerer_macros_v233.py à la génération v0.23.3 :
% les journaux vivants APRÈS la vague 15, la plage à la dernière décision,
% le DOI au jour de la rédaction ; les canons des vagues 9 à 15)
\newcommand{\nbLignesJournal}{13,651}
\newcommand{\nbLignesWorklog}{14,092}
\newcommand{\nbLignesDecisions}{4,848}
\newcommand{\plageDecisions}{D1--D269}
% LES DOI (le concept résout sur la dernière version)
\newcommand{\doiConcept}{10.5281/zenodo.24367343}
\newcommand{\doiVersionCourante}{10.5281/zenodo.24442951}
% LES CANONS M4 (le record daté de la vague 9 : n=5, la liste fixe, les
% médianes RECALCULÉES des JSON bruts — le ratio de régime courant
% RECALCULÉ des médianes non arrondies : 6.912/2.056 =
% 3.362 → 3.4 à une décimale, PARTOUT où le
% courant est cité)
\newcommand{\mquatreAvantMediane}{39.3}
\newcommand{\mquatreChaudMediane}{6.9}
\newcommand{\mquatreBrasC}{2.1}
\newcommand{\mquatreRatioAvant}{19.1$\times$}
\newcommand{\mquatreRatioChaud}{3.4$\times$}
% LE COMPTE D'ARBRES SIGNÉS (le sceau I — quatorze, D177)
\newcommand{\nbArbresSignes}{14}
% LES CANONS DE LA VAGUE 9 (reconduits — le record daté)
\newcommand{\sureteDecharges}{2006}
\newcommand{\sureteDechargesTotal}{2098}
\newcommand{\sureteNOMMES}{92}
\newcommand{\sureteMedies}{348}
\newcommand{\mutationBrutNum}{811}
\newcommand{\mutationDenomBrut}{832}
\newcommand{\mutationAjusteNum}{811}
\newcommand{\mutationDenomAjuste}{817}
\newcommand{\cfTroisObligations}{14/14}
% LES CANONS DE LA VAGUE 10 (reconduits — les pièces des ateliers)
\newcommand{\cgEffetNotationMs}{37073}
\newcommand{\cgEffetNotationRatio}{18.1$\times$}
\newcommand{\cgEffetGouvernanceMs}{215}
\newcommand{\cgEffetGouvernanceRatio}{1.11$\times$}
\newcommand{\cgMquatreT}{39.24}
\newcommand{\cgMquatreC}{1.95}
\newcommand{\cgMquatreCG}{2.16}
\newcommand{\fenetreEditerTotal}{1146}
\newcommand{\fenetreEditerReponse}{40}
\newcommand{\fenetreHdeuxTotal}{2452}
\newcommand{\fenetreReduction}{53.3}
\newcommand{\sseDeuxVdeuxInstructions}{4087}
\newcommand{\sseDeuxTempsReel}{833}
\newcommand{\gccODeuxTempsReel}{16040}
\newcommand{\idiomeMemchrTempsReel}{561}
\newcommand{\gardesNombre}{58}
\newcommand{\gardesLea}{34}
\newcommand{\niBlocsCouverts}{16}
\newcommand{\niFamilles}{7}
% LES CANONS DE LA VAGUE 11 (les pièces des ateliers, les médianes des
% JSON bruts)
% — LES QUATRE BRAS (3lmyfb2xa : la médiane de chaque bras, la
%   ventilation trois-effets RECALCULÉE : 39.393 − 18.678 =
%   20.715 ; la somme 20.715 + 16.497 +
%   0.213 = 37.425 ==
%   T−C 37.425)
\newcommand{\quatreBrasT}{39.39}
\newcommand{\quatreBrasCGP}{18.68}
\newcommand{\quatreBrasCG}{2.18}
\newcommand{\quatreBrasC}{1.97}
\newcommand{\effetNotationPropre}{20.72}
\newcommand{\effetRegime}{16.50}
\newcommand{\effetGouvernanceOnze}{0.213}
\newcommand{\partNotationPropre}{55.7}
\newcommand{\partRegime}{44.3}
% — LES COULOIRS SIMD (fnvan0zes : le rapport temps réel, les médianes)
\newcommand{\hybrideCycles}{704}
\newcommand{\hybrideIqr}{52}
\newcommand{\hybrideGain}{25.9$\times$}
\newcommand{\scelleGprCycles}{18226}
\newcommand{\shaniCorpsCbloc}{128.5}
\newcommand{\shaniGain}{6.5$\times$}
\newcommand{\avxDeuxCycles}{5846}
\newcommand{\cpuidSondage}{2{,}660}
\newcommand{\batterieVecteurs}{559}
% — LA BOUCLE PORTEUSE (ahult8jgd : la v4-inline, le cache, l'incrémental)
\newcommand{\porteuseVquatreInline}{7.09}
\newcommand{\porteuseVquatreIqr}{0.09}
\newcommand{\porteuseCible}{12}
\newcommand{\porteuseGain}{7.1$\times$}
\newcommand{\porteuseVdeuxFlux}{50.56}
\newcommand{\cacheRoundtripGain}{680.6$\times$}
\newcommand{\cristalIncremental}{54}
\newcommand{\oraclesPart}{81}
% — LE LIFTER AU DOMAINE (hrbdfy6i4)
\newcommand{\lifterEncodages}{942}
\newcommand{\lifterMutantesHorsVocab}{155}
% — LES GABARITS (o5apue194 : les médianes du champ prompt, le JSONL brut)
\newcommand{\gabaritEditerAvant}{1106}
\newcommand{\gabaritEditerApres}{674}
\newcommand{\gabaritEditerReduction}{39.1}
\newcommand{\gabaritLocaliserAvant}{1053}
\newcommand{\gabaritLocaliserApres}{627}
\newcommand{\gabaritLocaliserReduction}{40.5}
\newcommand{\gabaritAjouterAvant}{3074}
\newcommand{\gabaritAjouterApres}{2200}
\newcommand{\gabaritAjouterReduction}{28.4}
% — LE PLANCHER ÉDITER AU NATIF (zho5vilpy : la correction (c) — le
%   plancher re-tokenisé au champ usage du CLI z-ai, le MÊME instrument
%   que la campagne : médiane 27.5 natif ; le ratio
%   674/27.5 = 24.5)
\newcommand{\plancherEditerNatif}{27.5}
\newcommand{\ratioPlancherNatif}{24.5$\times$}
% — LA PURGE D81 (ou6wn6wlv)
\newcommand{\purgeCommitsScannes}{513}

% ===== LES CANONS v0.23.0 (ylrzmlbn8 : les opérandes non arrondis) =====
% le ratio ajouter : 1 − 2 199,5 / 3 073,5 = 28,43 % (jamais retapé)
\newcommand{\gabaritAjouterOp}{\% (opérandes 3\,073{,}5 $\to$ 2\,199{,}5)}
\newcommand{\gabaritLocaliserOp}{\% (opérandes 1\,053 $\to$ 627)}
\newcommand{\gabaritEditerOp}{\% (opérandes 1\,106 $\to$ 674)}
% LA PARTITION EXACTE DE LA BATTERIE 559 (la somme par code : 16+3+4+13+3+8+512)
\newcommand{\batteriePartition}{16 + 3 + 4 + 13 + 3 + 8 + 512}
% les opérandes non arrondis, la mise en forme ANGLAISE (les mêmes
% canons que les \gabarit*Op — l'abstract EN aligné, la tâche (d))
\newcommand{\gabaritLocaliserOpEN}{\% (operands 1{,}053 $\to$ 627)}
\newcommand{\gabaritEditerOpEN}{\% (operands 1{,}106 $\to$ 674)}
\newcommand{\gabaritAjouterOpEN}{\% (operands 3{,}073.5 $\to$ 2{,}199.5)}
